\documentclass[10pt,a4paper]{amsart}
\usepackage{amsmath,amssymb,amsthm,graphicx}
% Retained minimal preamble for the Stacks Project chapter more-algebra.tex (More on Algebra).
% The chapter file pinned-original-source.tex carries only \input{preamble} on its first line, so no
% preamble text lives inside the chapter itself.  The Stacks Project uses one single shared preamble for
% every chapter; the snapshot stages it identically as sources/stacks-chow/preamble.tex and
% sources/stacks-etale-cohomology/preamble.tex.  The block below is copied mechanically, line for line,
% from sources/stacks-chow/preamble.tex lines 14-18 (the Xy-pic setup) and 154-217 (the native theorem
% declarations together with the mathematical macros).  Nothing was added, removed or redefined; the
% wrapper class and the amsmath/amssymb/amsthm/graphicx base packages are supplied by the assembler.
% Chapter: more-algebra.tex -- More on Algebra.
% For commutative diagrams we use Xy-pic
\usepackage[all]{xy}

% We use 2cell for 2-commutative diagrams.
\xyoption{2cell}

% Theorem environments.
%
\theoremstyle{plain}
\newtheorem{theorem}[subsection]{Theorem}
\newtheorem{proposition}[subsection]{Proposition}
\newtheorem{lemma}[subsection]{Lemma}

\theoremstyle{definition}
\newtheorem{definition}[subsection]{Definition}
\newtheorem{example}[subsection]{Example}
\newtheorem{exercise}[subsection]{Exercise}
\newtheorem{situation}[subsection]{Situation}

\theoremstyle{remark}
\newtheorem{remark}[subsection]{Remark}
\newtheorem{remarks}[subsection]{Remarks}

\numberwithin{equation}{subsection}

% Macros
%
\def\lim{\mathop{\mathrm{lim}}\nolimits}
\def\colim{\mathop{\mathrm{colim}}\nolimits}
\def\Spec{\mathop{\mathrm{Spec}}}
\def\Hom{\mathop{\mathrm{Hom}}\nolimits}
\def\Ext{\mathop{\mathrm{Ext}}\nolimits}
\def\SheafHom{\mathop{\mathcal{H}\!\mathit{om}}\nolimits}
\def\SheafExt{\mathop{\mathcal{E}\!\mathit{xt}}\nolimits}
\def\Sch{\mathit{Sch}}
\def\Mor{\mathop{\mathrm{Mor}}\nolimits}
\def\Ob{\mathop{\mathrm{Ob}}\nolimits}
\def\Sh{\mathop{\mathit{Sh}}\nolimits}
\def\NL{\mathop{N\!L}\nolimits}
\def\CH{\mathop{\mathrm{CH}}\nolimits}
\def\proetale{{pro\text{-}\acute{e}tale}}
\def\etale{{\acute{e}tale}}
\def\QCoh{\mathit{QCoh}}
\def\Ker{\mathop{\mathrm{Ker}}}
\def\Im{\mathop{\mathrm{Im}}}
\def\Coker{\mathop{\mathrm{Coker}}}
\def\Coim{\mathop{\mathrm{Coim}}}

% Boxtimes
%
\DeclareMathSymbol{\boxtimes}{\mathbin}{AMSa}{"02}

%
% Macros for moduli stacks/spaces
%
\def\QCohstack{\mathcal{QC}\!\mathit{oh}}
\def\Cohstack{\mathcal{C}\!\mathit{oh}}
\def\Spacesstack{\mathcal{S}\!\mathit{paces}}
\def\Quotfunctor{\mathrm{Quot}}
\def\Hilbfunctor{\mathrm{Hilb}}
\def\Curvesstack{\mathcal{C}\!\mathit{urves}}
\def\Polarizedstack{\mathcal{P}\!\mathit{olarized}}
\def\Complexesstack{\mathcal{C}\!\mathit{omplexes}}
% \Pic is the operator that assigns to X its picard group, usage \Pic(X)
% \Picardstack_{X/B} denotes the Picard stack of X over B
% \Picardfunctor_{X/B} denotes the Picard functor of X over B
\def\Pic{\mathop{\mathrm{Pic}}\nolimits}
\def\Picardstack{\mathcal{P}\!\mathit{ic}}
\def\Picardfunctor{\mathrm{Pic}}
\def\Deformationcategory{\mathcal{D}\!\mathit{ef}}

% --- appended 2026-10-12 (Chief): the project own definitions of the reference/slogan/history
% display environments, copied mechanically from the shared Stacks preamble lines 8-12
% (\usepackage{verbatim} plus the three \newenvironment lines, which the chapter preamble keeps
% verbatim from the project preamble). Nothing added, removed or redefined: these environments
% carry no mathematical content and the project itself typesets them as comments. Without these
% three lines a counted statement whose own text uses them fails with Environment slogan undefined.
\usepackage{verbatim}
\newenvironment{reference}{\comment}{\endcomment}
\newenvironment{slogan}{\comment}{\endcomment}
\newenvironment{history}{\comment}{\endcomment}
\usepackage{cleveref}
\crefname{originallemma}{original lemma}{original lemmas}
\Crefname{originallemma}{Original lemma}{Original lemmas}
\makeatletter
\newlabel{algebra-lemma-characterize-etale}{{00U6}{source}{}{sourcepointer-1}{}}
\newlabel{algebra-lemma-characterize-etale@cref}{{[originallemma][0][]00U6}{[1][0][]source}{}{}{}}
\makeatother
\crefname{originallemma}{original lemma}{original lemmas}
\Crefname{originallemma}{Original lemma}{Original lemmas}
\makeatletter
\newlabel{algebra-lemma-etale}{{00U2}{source}{}{sourcepointer-2}{}}
\newlabel{algebra-lemma-etale@cref}{{[originallemma][0][]00U2}{[1][0][]source}{}{}{}}
\makeatother
\crefname{originallemma}{original lemma}{original lemmas}
\Crefname{originallemma}{Original lemma}{Original lemmas}
\makeatletter
\newlabel{fields-definition-normal-closure}{{0BMF}{source}{}{sourcepointer-3}{}}
\newlabel{fields-definition-normal-closure@cref}{{[originallemma][0][]0BMF}{[1][0][]source}{}{}{}}
\makeatother
\crefname{originallemma}{original lemma}{original lemmas}
\Crefname{originallemma}{Original lemma}{Original lemmas}
\makeatletter
\newlabel{fields-lemma-normal-closure-galois}{{0EXM}{source}{}{sourcepointer-4}{}}
\newlabel{fields-lemma-normal-closure-galois@cref}{{[originallemma][0][]0EXM}{[1][0][]source}{}{}{}}
\makeatother
\crefname{originallemma}{original lemma}{original lemmas}
\Crefname{originallemma}{Original lemma}{Original lemmas}
\makeatletter
\newlabel{fields-lemma-normal-closure-tensor-product}{{0EXL}{source}{}{sourcepointer-5}{}}
\newlabel{fields-lemma-normal-closure-tensor-product@cref}{{[originallemma][0][]0EXL}{[1][0][]source}{}{}{}}
\makeatother
\crefname{originallemma}{original lemma}{original lemmas}
\Crefname{originallemma}{Original lemma}{Original lemmas}
\makeatletter
\newlabel{lemma-Dedekind-etale-extension}{{10881}{source}{}{sourcepointer-6}{}}
\newlabel{lemma-Dedekind-etale-extension@cref}{{[originallemma][0][]10881}{[1][0][]source}{}{}{}}
\makeatother
\crefname{originallemma}{original lemma}{original lemmas}
\Crefname{originallemma}{Original lemma}{Original lemmas}
\makeatletter
\newlabel{lemma-multiplicative-e-f}{{32952}{source}{}{sourcepointer-7}{}}
\newlabel{lemma-multiplicative-e-f@cref}{{[originallemma][0][]32952}{[1][0][]source}{}{}{}}
\makeatother
\crefname{originalremark}{original remark}{original remarks}
\Crefname{originalremark}{Original remark}{Original remarks}
\makeatletter
\newlabel{remark-finite-separable-extension}{{33006}{source}{}{sourcepointer-8}{}}
\newlabel{remark-finite-separable-extension@cref}{{[originalremark][0][]33006}{[1][0][]source}{}{}{}}
\makeatother
\title{Unramified finite separable extensions of a discretely valued field are closed under subextension, Galois closure and compositum}
\author{The Stacks Project Authors}
\date{stacks-project:more-algebra.tex:local-snapshot-20261012; GNU FDL 1.2 or later}
\begin{document}
\maketitle
\setcounter{subsection}{1}
\noindent\textit{Editorial note.} All delivered mathematical text is the author's own, copied from the pinned Stacks Project chapter file more-algebra.tex; no statement or proof step was rewritten, reordered or removed.\par\medskip
\noindent\textit{Editorial note.} The chapter file pinned-original-source.tex carries only \textbackslash{}input\{preamble\} on its first line, so no preamble text lives inside the chapter. The retained file sources/stacks-more-algebra/retained-preamble.tex is a mechanical, line-for-line copy of the project's single shared preamble as staged in the snapshot; nothing was added, removed or redefined and no package was added.\par\medskip
\noindent\textit{Editorial note.} The source URL is the source-level Stacks Project URL, as required when more than one card is produced from one source; the per-claim tag is carried by claim\_key.\par\medskip
\noindent\textit{Editorial note.} The statement contains the author's duplicated article in 'there exists a a finite separable extension', and the proof ends with 'Details omitted.'; both are reproduced unchanged. The nine references inside the counted proof span are recorded as non-load-bearing navigation pointers.\par\medskip
\noindent\textit{Allowed standard prerequisites (editorial).}\begin{itemize}
\item The integral closure of a discrete valuation ring in a finite separable extension and multiplicativity of ramification indices. used at proof lines 33089-33096. Is the structure that lets the unramifiedness assumption be transported to a subextension.
\item Unramifiedness with respect to A as etaleness of A -> B. used at proof lines 33112-33115. Converts the hypothesis into a form to which permanence of etale ring maps applies.
\item The normal closure as a Galois extension receiving a surjection from a tensor power. used at proof lines 33102-33110. Produces the Galois extension required in part (2).
\item Permanence of etale ring maps under tensor products, and etale extensions of a discrete valuation ring as products of Dedekind domains. used at proof lines 33116-33123. Is the step showing the tensor power is a product of Dedekind domains, from which unramifiedness follows.
\end{itemize}\medskip
\ifdefined\hypertarget\hypertarget{sourcepointer-1}{}\fi
\noindent\textit{Source reference.} Original lemma 00U6: focused-seeds/AGP/sources/stacks-snapshot-20261011/tags.txt, source lines 1100--1100. Cross-chapter navigation pointer into the Stacks Project; the standard fact it supplies is applied inside the counted proof span and the argument resting on it stays there.\par\medskip
\ifdefined\hypertarget\hypertarget{sourcepointer-2}{}\fi
\noindent\textit{Source reference.} Original lemma 00U2: focused-seeds/AGP/sources/stacks-snapshot-20261011/tags.txt, source lines 1096--1096. Cross-chapter navigation pointer into the Stacks Project; the standard fact it supplies is applied inside the counted proof span and the argument resting on it stays there.\par\medskip
\ifdefined\hypertarget\hypertarget{sourcepointer-3}{}\fi
\noindent\textit{Source reference.} Original lemma 0BMF: focused-seeds/AGP/sources/stacks-snapshot-20261011/tags.txt, source lines 14628--14628. Cross-chapter navigation pointer into the Stacks Project; the standard fact it supplies is applied inside the counted proof span and the argument resting on it stays there.\par\medskip
\ifdefined\hypertarget\hypertarget{sourcepointer-4}{}\fi
\noindent\textit{Source reference.} Original lemma 0EXM: focused-seeds/AGP/sources/stacks-snapshot-20261011/tags.txt, source lines 18662--18662. Cross-chapter navigation pointer into the Stacks Project; the standard fact it supplies is applied inside the counted proof span and the argument resting on it stays there.\par\medskip
\ifdefined\hypertarget\hypertarget{sourcepointer-5}{}\fi
\noindent\textit{Source reference.} Original lemma 0EXL: focused-seeds/AGP/sources/stacks-snapshot-20261011/tags.txt, source lines 18661--18661. Cross-chapter navigation pointer into the Stacks Project; the standard fact it supplies is applied inside the counted proof span and the argument resting on it stays there.\par\medskip
\ifdefined\hypertarget\hypertarget{sourcepointer-6}{}\fi
\noindent\textit{Source reference.} Original lemma 10881: sources/stacks-more-algebra/pinned-original-source.tex, source lines 10881--10881. Same-chapter navigation pointer; the statement it names is standard for the section and the argument that consumes it stays inside the counted proof span.\par\medskip
\ifdefined\hypertarget\hypertarget{sourcepointer-7}{}\fi
\noindent\textit{Source reference.} Original lemma 32952: sources/stacks-more-algebra/pinned-original-source.tex, source lines 32952--32952. Same-chapter navigation pointer; the statement it names is standard for the section and the argument that consumes it stays inside the counted proof span.\par\medskip
\ifdefined\hypertarget\hypertarget{sourcepointer-8}{}\fi
\noindent\textit{Source reference.} Original remark 33006: sources/stacks-more-algebra/pinned-original-source.tex, source lines 33006--33006. Same-chapter navigation pointer; the statement it names is standard for the section and the argument that consumes it stays inside the counted proof span.\par\medskip
\begin{lemma}
\label{lemma-permanence-unramified}
Let $A$ be a discrete valuation ring with fraction field $K$.
\begin{enumerate}
\item If $M/L/K$ are finite separable extensions and
$M$ is unramified with respect to $A$, then $L$ is unramified
with respect to $A$.
\item If $L/K$ is a finite separable extension which is
unramified with respect to $A$, then there exists a Galois
extension $M/K$ containing $L$ which is unramified with respect to $A$.
\item If $L_1/K$, $L_2/K$ are finite separable extensions which are
unramified with respect to $A$, then there exists a a finite
separable extension $L/K$ which is unramified with respect
to $A$ containing $L_1$ and $L_2$.
\end{enumerate}
\end{lemma}

\begin{proof}
We will use the results of the discussion in
Remark \ref{remark-finite-separable-extension}
without further mention.

\medskip\noindent
Proof of (1). Let $C/B/A$ be the integral closures of $A$ in $M/L/K$.
Since $C$ is a finite ring extension of $B$, we see that
$\Spec(C) \to \Spec(B)$ is surjective. Hence for ever maximal
ideal $\mathfrak m \subset B$ there is a maximal ideal
$\mathfrak m' \subset C$ lying over $\mathfrak m$.
By the multiplicativity of ramification indices
(Lemma \ref{lemma-multiplicative-e-f})
and the assumption, we conclude that the ramification
index of $B_\mathfrak m$ over $A$ is $1$.
Since $\kappa(\mathfrak m')/\kappa_A$ is finite separable,
the same is true for $\kappa(\mathfrak m)/\kappa_A$.

\medskip\noindent
Proof of (2). Let $M$ be the normal closure of $L$ over $K$, see
Fields, Definition \ref{fields-definition-normal-closure}.
Then $M/K$ is Galois by Fields, Lemma \ref{fields-lemma-normal-closure-galois}.
On the other hand, there is a surjection
$$
L \otimes_K \ldots \otimes_K L \longrightarrow M
$$
of $K$-algebras, see Fields, Lemma
\ref{fields-lemma-normal-closure-tensor-product}.
Let $B$ be the integral closure of $A$ in $L$
as in Remark \ref{remark-finite-separable-extension}. The
condition that $L$ is unramified with respect to $A$
exactly means that $A \to B$ is an \'etale ring map, see
Algebra, Lemma \ref{algebra-lemma-characterize-etale}.
By permanence properties of \'etale ring maps we see that
$$
B \otimes_A \ldots \otimes_A B
$$
is \'etale over $A$, see
Algebra, Lemma \ref{algebra-lemma-etale}.
Hence the displayed ring is a product of Dedekind domains, see
Lemma \ref{lemma-Dedekind-etale-extension}.
We conclude that $M$ is the fraction field of a
Dedekind domain finite \'etale over $A$.
This means that $M$ is unramified with respect to $A$
as desired.

\medskip\noindent
Proof of (3). Let $B_i \subset L_i$ be the integral closure of $A$.
Argue in the same manner as above to show that $B_1 \otimes_A B_2$
is finite \'etale over $A$. Details omitted.
\end{proof}

\end{document}
